\documentclass[../main.tex]{subfiles} % 使用 subfiles 文档类，指定主文档
\graphicspath{{\subfix{../image/}}} % 指定图片目录，后续可以直接使用图片文件名。

% 例如:
% \begin{figure}[h]
% \centering
% \includegraphics{image-04.04}
% \label{fig:image-04.04}
% \caption{图片标题}
% \end{figure}

\begin{document}

\chapter{波动方程}

\begin{definition}[波动方程]
设 $\Omega\subset\mathbb{R}^n$，$a>0$ 为常数，$f=f(x,t)$ 为已知函数，$u=u(x,t)$，$x\in\Omega$，$t>0$. 称方程
\begin{align}
u_{tt}-a^2\Delta u=f \label{eq:wave-gen}
\end{align}
为\textbf{波动方程}. 当 $f\equiv 0$ 时，称
\begin{align*}
u_{tt}-a^2\Delta u=0
\end{align*}
为\textbf{齐次波动方程}.
\end{definition}

\begin{definition}
设 $Q\subset \mathbb{R}^n\times(0,\infty)$，我们用 $C^{2,2}(Q)$ 表示所有在 $Q$ 内关于 $x$ 二阶偏导数连续、关于 $t$ 二阶偏导数连续的函数构成的集合，即
\begin{align*}
C^{2,2}(Q) \triangleq \{\, u \mid u, u_t, u_{tt}, u_{x_i}, u_{x_ix_j} \in C(Q),\ i,j=1,\cdots,n\,\}.
\end{align*}
称波动方程 \eqref{eq:wave-gen} 在 $C^{2,2}$ 函数集中的解为波动方程的\textbf{古典解}.
\end{definition}

\begin{definition}[波动方程的定解条件]
波动方程 \eqref{eq:wave-gen} 的定解条件包括初始条件和边值条件：
\begin{enumerate}[(1)]
\item \textbf{初始条件：} 给出弹性体各点在初始时刻 $t=0$ 的位移和速度：
\begin{align}
\begin{cases}
u(x,0)=\varphi(x), & x\in\Omega, \\
u_t(x,0)=\psi(x), & x\in\Omega,
\end{cases} \label{eq:wave-init}
\end{align}
其中 $\varphi,\psi$ 为已知函数. 称
\begin{align}
\begin{cases}
u_{tt}-a^2\Delta u=f(x,t), & (x,t)\in\mathbb{R}^n\times\mathbb{R}_+,\\
u(x,0)=\varphi(x),\quad u_t(x,0)=\psi(x), & x\in\mathbb{R}^n,
\end{cases}
\end{align}
为\textbf{$n$维波动方程初值问题（Cauchy初值问题）}.

\item \textbf{边值条件：} 通常有以下三类：
\begin{enumerate}[(i)]
\item \textbf{第一类边值条件（Dirichlet 条件）：}
\begin{align}
u(x,t)=g(x,t),\qquad x\in\partial\Omega,\ t>0. \label{eq:wave-bc1}
\end{align}
当 $g\equiv 0$ 时，表示边界固定.

\item \textbf{第二类边值条件（Neumann 条件）：}
\begin{align}
\frac{\partial}{\partial n}u(x,t)=g(x,t),\qquad x\in\partial\Omega,\ t>0, \label{eq:wave-bc2}
\end{align}
其中 $n$ 是 $\partial\Omega$ 的单位外法向量. 当 $g\equiv 0$ 时，表示边界自由（无外力作用）.

\item \textbf{第三类边值条件（Robin 条件）：}
\begin{align}
\frac{\partial}{\partial n}u(x,t)+\alpha(x,t)u(x,t)=g(x,t),\qquad x\in\partial\Omega,\ t>0, \label{eq:wave-bc3}
\end{align}
其中 $\alpha(x,t)>0$，$n$ 是 $\partial\Omega$ 的单位外法向量. 当 $g\equiv 0$ 时，表示边界固定在弹性支撑上.
\end{enumerate}
\end{enumerate}
称波动方程 \eqref{eq:wave-gen} 连同初始条件 \eqref{eq:wave-init} 及某一类边值条件构成的问题为波动方程的\textbf{混合问题}（或初边值问题）.
\end{definition}
\begin{remark}
波动方程 \eqref{eq:wave-gen} 的物理背景：若 $u(x,t)$ 表示物体在位置 $x$、时刻 $t$ 的位移，则 $a=\sqrt{T_0/\rho_0}$，其中 $T_0$ 为弹性内力强度，$\rho_0$ 为密度，$f$ 为外力密度除以密度. 当 $n=1$ 时，描述弦的振动；$n=2$ 时描述薄膜振动；$n=3$ 时描述弹性体振动.
\end{remark}

\subfile{Chapter-04/section-01.tex}

\subfile{Chapter-04/section-02.tex}

\subfile{Chapter-04/section-03.tex}

\subfile{Chapter-04/section-04.tex}

\subfile{Chapter-04/section-05.tex}

\subfile{Chapter-04/section-06.tex}

\subfile{Chapter-04/section-07.tex}

\subfile{Chapter-04/section-08.tex}

\subfile{Chapter-04/section-09.tex}

\subfile{Chapter-04/section-10.tex}

\end{document}